\phantomsection
\subsection*{M1A. Dry Mead}
\addcontentsline{toc}{subsection}{M1A. Dry Mead}

\textbf{Impressão Geral}: É um hidromel sem dulçor residual significativo, permitindo aromas naturais do mel a serem apreciados sem dulçor associado a mel cru. Baixos níveis de dulçor podem ser percebidos devido a fermentação, maturação ou envelhecimento. Possui equilíbrio, corpo, retrogosto e intensidade de sabor similar ao encontrado em diversos exemplares de vinhos brancos secos.

\textbf{Aroma}: Sutil a perceptível aroma de mel, reflexo de méis varietais declarados. Aromas significativos de mel ou impressão de dulçor não são esperados. Os aromas do mel são mais percebidos na forma de aromas florais, frutados, condimentados ou outras notas que venham do mel, ao invés de aromas doces de méis crus não fermentados. Os baixos níveis de dulçor permitem que os aromas de fermentação e o caráter de envelhecimento sejam perceptíveis. Hidroméis mais alcoólicos podem apresentar leves notas condimentadas oriundas do álcool. Hidroméis carbonatados tendem a expressar mais aromas.

\textbf{Aparência}: De limpida a brilhante. Deseja-se uma maior limpidez. Quaisquer bolhas, espuma ou efervescência são baseadas no nível de carbonatação declarado. A cor é proveniente das variedades de méis declarados, mas geralmente de amarelo palha a dourado. Cores mais vivas e brilhantes são mais desejadas.

\textbf{Sabor}: Sutil a perceptível caráter de mel, reflexo dos méis varietais declarados. Sabores semelhantes aos aromas de mel fermentado em Aroma. Final seco, de suave a crocante, mas não totalmente desprovido de dulçor. Níveis de dulçor residual são mínimos a nenhum. Acidez pode ser percebida, mas não deve ser afiada ou cortante. Taninos podem tornar o hidromel ainda mais seco. Características ásperas, desagradáveis ou excessivamente sulfurosas ou com notas de levedura/autólise são indesejáveis.

\textbf{Sensação na Boca}: Corpo geralmente de baixo a médio, mas não deve parecer aguado. Hidroméis com maior força alcoólica podem apresentar corpo mais cheio e maior aquecimento alcoólico. Sensação de corpo não deve ser acompanhada de dulçor residual perceptível. A carbonatação é indicada pelo nível declarado, de tranquilo a espumante.

\textbf{Ingredientes}: Mel de qualquer origem. Espera-se que os méis varietais apresentem características distintas associadas às variedades declaradas.

\textbf{Comentários}: Podem parecer com vinhos mais do que outros tipos de hidromel uma vez que o dulçor associado ao mel é amplamente ausente. Hidroméis secos podem possuir mais ênfase nas características residuais de fermentação do mel e da levedura, uma vez que esses são os ingredientes principais que proverão complexidade.

\textbf{Instruções para Inscrição}: Participantes devem declarar dulçor, carbonatação e força alcoólica; dulçor é restrito a seco nesse estilo. Variedades de mel podem ser declaradas.

\textbf{Exemplos Comerciais}: Brothers Drake Wild Ohio, Honeygirl Meadery Orange Blossom, Heidrun Meadery (todos exemplares de espumantes secos de méis monoflorais), Intermiel Bouquet Printanier, White Winter Dry Mead